\subsection{A gonality bound for $X_1(2,32)$}
\begin{lemma}\label{lem:decic232-gonality}
The $\Q$-gonality of $X_1(2,32)$ is at least $12$.
\end{lemma}
\begin{proof}
Put
\[
 \Delta_0=\{\pm1,\pm31\},\qquad
 \Delta_1=\{\pm1,\pm15,\pm17,\pm31\}
 \subseteq(\mathbb Z/64\mathbb Z)^*.
\]
The usual $2$-isogeny construction identifies $X=X_1(2,32)$
with $X_{\Delta_0}(64)$ over $\Q$.
Indeed, from $(E,P,Q)$ take $\phi:E\to E'=E/\langle P\rangle$
and a point $R$ with $\widehat\phi(R)=Q$. Then
$2R=\phi(Q)$ has order $32$, so $R$ has order $64$.
The two choices of $R$ differ by $32R$ and therefore give the
same $\Delta_0$-orbit. Conversely the dual isogeny recovers the
original level structure. This also follows by conjugating
$\Gamma_1(32)\cap\Gamma^0(2)$ by $\operatorname{diag}(1,2)$.

The inclusion $\Delta_0\subset\Delta_1$ gives a degree-two map
\[
 \pi:X\longrightarrow Y=X_{\Delta_1}(64).
\]
Here $g(X)=37$ and $g(Y)=13$. By
\cite[Proposition~3.7 and Table~2]{Orlic},
$\operatorname{gon}_{\Q}Y\geq6$.
Suppose that $f:X\to\mathbb P^1$ has degree $d\leq11$.
If $f$ factors through $\pi$, its induced map on $Y$ has degree
$d/2\leq5$, a contradiction. Otherwise the two maps generate
$\Q(X)$, since $[\Q(X):\Q(Y)]=2$.
Castelnuovo--Severi gives
\[
 37=g(X)\leq2g(Y)+(2-1)(d-1)=25+d\leq36,
\]
again a contradiction.
\end{proof}


\begin{lemma}\label{lem:decic232-index-two}
Let $Y/\Q$ be a curve with a rational cusp $c_0$ and finite
$J_Y(\Q)$. Let $H_\Q$ be the subgroup generated by its rational
cuspidal divisor classes, and suppose that $[J_Y(\Q):H_\Q]\leq2$.
If $g$ is a rational automorphism preserving the cusps, then, for
every rational divisor $B$ of degree $d$,
\[
 [gB-dc_0]\equiv[B-dc_0]\pmod{H_\Q}.
\]
\end{lemma}
\begin{proof}
The group $J_Y(\Q)/H_\Q$ has order at most two, so its automorphism
induced by $g$ is the identity. Moreover $[gc_0-c_0]\in H_\Q$.
The assertion follows by applying $g$ to $[B-dc_0]$.
\end{proof}

\begin{proposition}\label{prop:decic232}
No elliptic curve over a decic number field contains
$C_2\times C_{32}$ or $C_{64}$.
\end{proposition}
\begin{proof}
Put $X=X_1(2,32)$. Its genus is $37$, its Jacobian has rank zero
by \cite[Theorem~3.1]{DEHMZ}, and it has good reduction at $3$.
Suppose that a decic level point exists, and let $D$ be its
noncuspidal trace divisor of degree $10$.
Lemma~\ref{lem:decic232-gonality} and
Lemma~\ref{lem:decic-cuspidal-reduction} exclude entirely cuspidal
reduction. Thus write $\overline D=N+C$, where $N$ has
noncuspidal support and $C$ is cuspidal.

The complete noncuspidal closed-point counts over $\mathbb F_3$,
in degrees one through ten, are
\[
 (0,0,0,0,0,192,320,960,1824,6208).
\]
They agree with Newton identities for the Frobenius polynomial of
$J_H(64)$ with $H=\langle33\rangle$.
The cusp signature is $1^{16},2^8,4^2,8^2$.
Consequently $N$ is a single closed point of degree $6$, $7$,
$8$, $9$ or $10$. The rational level automorphisms
\[
 (P,Q)\longmapsto(P+16bQ,cP+aQ),\quad b,c\in\{0,1\},\quad
 a\in\{1,3,5,7,9,11,13,15\},
\]
reduce these cases to respectively $8,10,32,57,198$ orbits.
Every one of the $9504$ closed points is accounted for.

Apply the four maps to $X_1(32)$ that retain $Q$ or $P+Q$, or
retain the image of $Q$ after quotienting by $\langle P\rangle$
or $\langle P+16Q\rangle$.
By \cite[Theorem~4.13 and Table~2]{DEHMZ}, the rational Jacobian
of $X_1(32)$ is generated by rational cuspidal divisor classes and
has invariants $[2,2,2,4,120,11640]$.
At $3$ the Jacobian has invariants
$[2,2,2,4,4,120,197880]$, so the rational image has index $68$.
Allowing every finite cuspidal complement gives the following sieve:
\[
\begin{array}{c|rrrrr}
\deg N&6&7&8&9&10\\\hline
\text{level orbits}&8&10&32&57&198\\
\text{surviving orbits}&1&0&0&0&0.
\end{array}
\]
Image residue degrees and their norm multiplicities are included.
All finite cusps are allowed, whether or not their divisors lift
rationally. Independent finite-field arithmetic and exact ideal
identities verify the sieve.

We describe the remaining orbit. Over
\[
 k=\mathbb F_3[\alpha]/
 (\alpha^6+2\alpha^4+\alpha^2+2\alpha+2)
\]
it has the normalized representative
\[
 E:y^2=x^3+cx^2+b(1-b-c)x,\quad P=(0,0),\quad Q=(b,b),
\]
where
\[
 b=\alpha^5,\qquad
 c=1+\alpha^2+\alpha^3+2\alpha^4+\alpha^5.
\]
For each of the four maps above its class admits a degree-four
cuspidal complement, but no complement of degree at most three.
Every $\mathbb F_3$-rational source cusp is globally rational.
If $C$ contained such a cusp, removing its image would leave the
same target class modulo the rational image and a complement of
degree three. Hence $C$ has no rational cusp in its support.

Put $h(P,Q)=(P,15Q)$. On the displayed point, $h$ acts as the
third power of Frobenius, so $hN=N$.
The Tate cusp calculation shows that $h$ fixes geometrically every
cusp in a closed point of degree two or four. Thus $hC=C$ and
$h\overline D=\overline D$.
For clarity, one can check these statements using primitive labels
\[
 Q=\zeta_{32}^a q^{r/32},\qquad P=(-1)^s q^{t/2},
 \qquad sr-ta\equiv1\pmod2,
\]
modulo unipotent translation and simultaneous negation.
Galois acts by $a\mapsto ua$, and $h$ acts by
$(a,r)\mapsto(15a,15r)$. This gives the stated cusp signature,
the sixteen globally rational cusps, and the required fixed cusps.

The class $[hD-D]$ is rational torsion and reduces to zero at $3$.
Torsion injectivity at this unramified odd prime gives $hD\sim D$.
Since $\deg D=10<\operatorname{gon}_{\Q}X$, we have $hD=D$.
The involution has no noncuspidal geometric fixed point in
characteristic zero: an elliptic-curve automorphism realizing $h$
would have square one, whereas multiplication by $15$ on the
point of order $32$ is neither $1$ nor $-1$.
It follows that
\[
 D=\pi^*B,\qquad \deg B=5,
 \qquad \pi:X\longrightarrow Y=X/\langle h\rangle.
\]
Here $Y=X_{\Delta_1}(64)$ from
Lemma~\ref{lem:decic232-gonality}, so $g(Y)=13$.
The cubic image $N_3=\pi(N)$ is noncuspidal.
The cusps in $C$ are ramified in $\pi$, so their multiplicities
in $\overline D$ are even. Consequently
\[
 \overline B=N_3+C_2
\]
for one of the eight quadratic cusps of $Y_{\mathbb F_3}$.
This argument applies to effective trace divisors as well as
primitive degree-ten points.

We now work on $Y$. Its Jacobian has rank zero, and its proven
finite Jacobian and rational cuspidal image are
\[
\begin{aligned}
 G=J_Y(\mathbb F_3)&=[2,2,2,2,8,40,40,80],\\
 H&=[2,2,4,20,40,40].
\end{aligned}
\]
Thus $|H|=512000$ and $G/H=C_2\times C_2\times C_8$.
The twenty finite cusps form seventeen global Galois orbits.
Their orbit sums give the stated subgroup $H$.
An exact modular-symbol calculation, using the full integral
homology lattice, gives
\[
 [J_Y(\Q):H_\Q]\mid2.
\]
Here are the two parts of this certificate. The operators
$T_3-\langle3\rangle-3$, $T_5-\langle5\rangle-5$, and complex
conjugation minus one give the Hecke bound in the $X_\mu$
convention. Fricke identifies this convention with the
$\Q$-isomorphic $X_{\Delta_1}(64)$ model.
The quotient of this bound by its intersection with the geometric
cuspidal subgroup has order two. The Galois-fixed geometric
cuspidal classes are exactly the classes of rational cuspidal
divisors, of order $512000$.
These are computations on integral lattices of rank $26$;
no odd-primary restriction is imposed.
The exported period, Hecke, conjugation and Galois lattices verify
both assertions. This is the Hecke-bound method of
\cite[Section~4.2]{DEHMZ}.

The diamond $\langle3\rangle$ commutes with $h$ and descends to $Y$.
By Lemma~\ref{lem:decic232-index-two}, the reductions of
$[B-5c_0]$ and $[\langle3\rangle B-5c_0]$ must have the same
class modulo $H$, and this class has order dividing two.
For $g=1,\langle3\rangle$, form the necessary sets
\[
 S_g=\{[gN_3+C_2-5c_0]\bmod H:
 C_2\text{ a quadratic cusp},\quad
 2[gN_3+C_2-5c_0]\in H\}.
\]
In a single exported basis of $G/H=C_2\times C_2\times C_8$, they are
\[
 S_1=\{(1,1,0),(1,0,4)\},\qquad
 S_{\langle3\rangle}=\{(1,1,4),(1,0,0)\}.
\]
They are disjoint, a contradiction.
All these classes are direct lookups of actual places in the
proven finite Picard presentation. A fresh computation proves
the presentation again and identifies its exact relation kernel;
independent finite-field and integer-lattice arithmetic verifies
the moduli maps and the disjoint sets.
Thus $C_2\times C_{32}$ is excluded.
Lemma~\ref{lem:decic-isogeny-transfer} excludes $C_{64}$ as well.
\end{proof}
